% claim-ref: java-moss#method
% claim-ref: java-moss#contain-discard
% claim-ref: java-moss#glue-product
% claim-ref: java-moss#glue-hold
% claim-ref: java-moss#glue-contact
% claim-ref: java-moss#glue-safety
% claim-ref: java-moss#observe-success
\CompanionHeader{PROPAGATION}{Java moss}{Taxiphyllum barbieri}{Aquatic plant / moss portions attached to decor. Proposed home sequence; identity and tank readings pending.}
\Context{Start only with healthy material and a prepared destination. Keep aquatic tissue wet; do not import terrestrial soil, callusing or watering schedules.}
\begin{minipage}[t]{.625\linewidth}\raggedright\vspace{0pt}
% step-claims: java-moss#contain-discard
\Step{1}{Prepare / inspect}{Use clean scissors, tank water, removable wood or stone and a waste bag. Inspect new material before it enters the aquarium. \textbf{[DEC]}}
% step-claims: java-moss#method
\Step{2}{Select / divide}{Separate a healthy small moss portion; retain growing tips. No minimum portion size is established. \textbf{[JM-TROP, JM-COOP]}}
% step-claims: java-moss#method java-moss#glue-product
\Step{3}{Choose attachment}{Tie gently to decor with thread/line, or use the optional aquarium-labeled cyanoacrylate gel method below. Check ties for loose loops. \textbf{[JM-TROP, GLUE]}}
% step-claims: java-moss#glue-contact
\Step{4}{Apply gel sparingly}{For Flourish Glue, use a small amount at moss contact points on wood or stone. Leave most of each portion and its growing tips uncoated. Do not glue to glass. \textbf{[GLUE]}}
% step-claims: java-moss#glue-hold
\Step{5}{Hold / inspect}{Manufacturer instruction: press to the surface for 20 seconds. This is a hold instruction, not proof of complete cure or biological establishment. Follow the actual label. \textbf{[GLUE]}}
% step-claims: java-moss#method
\Step{6}{Place / observe}{Keep the moss submerged and loose enough to inspect; avoid smothering it in a thick glued clump. \textbf{[JM-TROP]}}
% step-claims: java-moss#contain-discard java-moss#observe-success
\Step{7}{Maintain / contain}{Observe continued growth and eventual natural hold without pulling. Trim excess, capture pieces and seal unwanted fragments for trash; no outdoor release or compost. \textbf{[DEC, JM-TROP]}}
\end{minipage}\hfill\begin{minipage}[t]{.345\linewidth}\raggedright\vspace{0pt}
\Note{Milestones}{Separate material; observe new growth; confirm continued healthy growth. A cut or glue bond is only the starting event. No universal days-to-success is established.}
\Note{Product-specific glue}{Flourish Glue is cyanoacrylate gel labeled for aquarium use, including underwater. Do not assume an unknown household glue has the same formulation or instructions. \textbf{[GLUE]}}
\Note{Handling}{Read the label. Cyanoacrylate bonds skin and eyes rapidly and is an irritant. Keep away from children. Avoid skin/eye contact. \textbf{[GLUE]}}
\Note{If growth stalls}{Record light, water, flow and recent changes. Inspect for detached, buried or deteriorating tissue. These are checks, not a diagnosis or medication instruction.}
\Note{No safety guarantee}{Plants do not prove cycling, sufficient oxygen or predation-safe shelter. Review animal care and current water results separately.}
\end{minipage}\par\vspace{4pt}
\Panel{Observation record}{Start / method: \hrulefill\quad Date: \hrulefill\par New growth / attachment / conditions: \hrulefill}
\Sources{\href{https://tropica.com/en/plants/plantdetails/4391/4391}{JM-TROP: Tropica Aquarium Plants}; \href{https://www.aquariumcoop.com/blogs/aquarium/java-moss}{JM-COOP: Aquarium Co-Op}; \href{https://www.seachem.com/flourish-glue.php}{GLUE: Seachem Laboratories}; \href{https://dec.ny.gov/nature/animals-fish-plants/invasive-species/aquatic/prevent-spread-of-aquatic-invasive-species/guidelines-for-aquarium-and-pet-owners}{DEC: New York State Department of Environmental Conservation}. Full scopes and claim records: adjacent YAML. Proposed instructions adapt cited methods; local conditions remain unknown.}{java-moss / propagation}
